\documentclass[11pt]{article}
\usepackage[T1]{fontenc}
\usepackage[utf8]{inputenc}
\usepackage{amsmath,amssymb,amsthm}
\usepackage[a4paper,margin=1in]{geometry}
\usepackage{booktabs,tabularx}
\usepackage[hidelinks]{hyperref}
\setlength{\emergencystretch}{3em}

\newtheorem{proposition}{Proposition}
\newcommand{\choosef}[2]{\binom{#1}{#2}}
\newcolumntype{Y}{>{\raggedright\arraybackslash}X}

\title{Erd\H{o}s Problem 700(iii):\\
Research Map and Open Frontier}
\author{Nanocodex Erd\H{o}s 700 project}
\date{3 August 2026}

\begin{document}
\maketitle

\begin{abstract}
This is a status report, not a claimed solution.  It records the exact
weighted carry reduction for Erd\H{o}s Problem 700(iii), the strongest
retained partial results, the approaches for which we found rigorous
counterexamples or quantifier failures, and the remaining same-row
synchronization problem.  No work in this repository proves the target for
any fixed exponent $A>1$, gives an unbounded improvement over the classical
$A=1$ scale for all composite integers, or constructs a counterexample to the
original statement.
\end{abstract}

\section{Target and exact reformulation}

For composite $n>1$, define
\[
 f(n)=\min_{2\leq k\leq\lfloor n/2\rfloor}
       \gcd\!\left(n,\choosef nk\right).
\]
Part (iii) asks whether, for every real $A>0$, there is a constant $C_A>0$
such that
\[
 f(n)\leq C_A\frac{n}{(\log n)^A}                         \tag{1}
\]
for every composite $n$.  Put
\[
 D_n(k)=\frac{n}{\gcd(n,\choosef nk)},\qquad
 D(n)=\max_{2\leq k\leq\lfloor n/2\rfloor}D_n(k).
\]
Then (1) is equivalent, after absorbing finitely many small $n$, to
\[
 D(n)\geq c_A(\log n)^A                                   \tag{2}
\]
for a constant $c_A>0$ depending only on $A$.

Write $n=\prod_{p\mid n}p^{a_p}$ and define the retained depth at a legal row
by
\[
 d_p(k)=\max\!\left\{a_p-v_p\!\left(\choosef nk\right),0\right\}.
\]
Primewise valuation gives the exact objective
\[
 \boxed{
 \log D_n(k)=\sum_{p\mid n}d_p(k)\log p.}                  \tag{3}
\]
Kummer identifies the binomial valuation with the number of base-$p$ carries.
Thus the open problem is to find one literal legal row $k$ retaining total
weighted depth at least
\[
 A\log\log n-O_A(1).                                      \tag{4}
\]
The words ``one row'' are load-bearing: successes at different primes or
scales do not combine unless they occur at the same integer $k$.

If $q=p^a\parallel n$ and $q\mid k$, scaling gives
\[
 v_p\!\left(\choosef nk\right)
 =v_p\!\left(\choosef{n/q}{k/q}\right),
 \qquad
 d_p(k)=\max\!\left\{a-v_p\!\left(\choosef{n/q}{k/q}\right),0\right\}.
                                                                    \tag{5}
\]
This exact component language underlies all of the retained routes.

\section{What has been proved}

The following are complete partial results, not substitutes for (1).

\begin{enumerate}
\item \textbf{Classical and easy regimes.}
The greatest-component witness and the classical least-common-multiple
estimate recover the $A\leq1$ scale.  Prime powers, boundedly many exact
prime-power components, a sufficiently large component, and sufficiently
large high-height mass are also controlled.

\item \textbf{Exact hard-case reduction.}
Outside those regimes, the obstruction can be reduced to many comparable,
low-height exact prime-power components.  Ordering and pigeonholing produce a
homogeneous dyadic packet of growing prime bases.  The component cap,
high-height cutoff, integral-$A$ endpoint, and legal-row endpoint were checked
explicitly.

\item \textbf{Deep-prefix escape.}
For each selected fixed-size subset, explicit Chinese-remainder rows can
suppress carries through
$\Omega_A(\log n/\log\log n)$ digit levels.  Therefore a genuine obstruction
must recur in the deep tail; bounded-prefix counterexamples cannot settle the
problem.

\item \textbf{Exact effective conductor.}
For
\[
 \mathcal A=\{0\leq z<p^D:z\leq T, v_p\choosef Tz<s\},
 \qquad \Delta=p^D-T,
\]
the least residue conductor retained by the campaign is
\[
 E=
 \begin{cases}
 0,&\Delta=1,\\
 \min\!\left(D,\lceil\log_p\Delta\rceil+s-1\right),&\Delta>1.
 \end{cases}                                                \tag{6}
\]
This replaces an abstract tail period by an exact top-gap statistic.

\item \textbf{Limits on structured blockers.}
Exact power-word anchors in a short prime packet have pairwise-coprime
exponents, yielding a sub-packet upper bound on their number.  A separate
cyclotomic argument excludes one sparse four-digit endpoint architecture for
fixed $A<2$.  Neither result classifies all hostile digit words.

\item \textbf{Conditional transfer.}
Any construction of one legal row satisfying (4) proves Part (iii).  The
often-used upper-half first-layer condition is a stronger sufficient
condition, not an equivalent reformulation.
\end{enumerate}

\section{Attempt map and no-retry lessons}

\begin{center}
\small
\begin{tabularx}{\textwidth}{@{}>{\raggedright\arraybackslash}p{0.24\textwidth}Y@{}}
\toprule
Route family & Disposition and retained lesson \\
\midrule
Fixed multiplier menus and naive witness iteration & Explicit counterexamples show that no fixed finite menu, copied component row, or local bounded-loss peeling rule preserves enough depth uniformly. \\
Independent carry probabilities and abstract set cover & Marginal success at separate primes does not produce one common row.  Any averaging proof needs arithmetic correlation in the actual signed phases. \\
Pair and packet alignment & Several universal pair-selection and upper-half full-layer statements are false.  An instance-adaptive multiplier or an inverse theorem for hostile words remains plausible. \\
Carmichael, lcm, determinant, and resultant bridges & Useful exact identities survive, but a fixed cubic Carmichael bridge is false and bounded-degree algebra loses the growing digit tail. \\
Digit boxes and affine rigidity & Exact escape bounds and sharp examples were proved.  Total box entropy or active-digit count alone cannot force a positive fraction of useful multipliers. \\
Large sieve, characters, and moment methods & Taking absolute values destroys the phase needed for simultaneous moving-base control.  Fixed moments stop at the singleton/$A=1$ scale. \\
Counterexample construction & We found counterexamples to methods, not to Erd\H{o}s 700(iii).  No infinite family with a uniform all-row bound contradicting (2) is known here. \\
\bottomrule
\end{tabularx}
\end{center}

These routes should not be restarted unchanged.  A new version must state
the additional arithmetic input that defeats the recorded counterexample or
repairs the failed quantifier.

\section{The surviving bottleneck}

The most direct target is a weighted-depth preservation theorem: move inside
an explicit legal progression so as to gain a new prime-power layer while
losing only a controlled fraction of previously retained weight.  A finite
iteration would need to reach (4).

For $1<A\leq2$, the residual packet is essentially squarefree, so extra
prime-power height gives no slack.  The central pair problem is therefore:

\begin{quote}
Given the actual common cofactor and a sufficiently large packet of primes,
find distinct $p,q$ and one legal multiplier satisfying both complete
base-$p$ and base-$q$ Lucas digit inequalities.
\end{quote}

A plausible route is an inverse theorem: long affine avoidance of a Lucas
digit down-set forces the upper word into a sparse or near-power class; a
moving-base separation theorem would then show that only a negligible
fraction of packet primes can be hostile for the same ambient integer.
Alternatively, an all-row contradiction could average the exact quantities
in (3), but it must retain their signed cross-base phase instead of multiplying
marginal probabilities.

\section{What would count as progress}

The repository's acceptance gate for Part (iii) is intentionally concrete.
Any one of the following would materially move the frontier:
\begin{enumerate}
\item a uniform lower bound $D(n)\geq L(n)\log n$ with $L(n)\to\infty$;
\item a proof of (2) for $A=2$;
\item the full theorem for every fixed $A>0$;
\item an infinite counterfamily negating (2) for one fixed $A>0$;
\item a carry-preserving composition or moving-base inverse theorem strong
enough to imply (4).
\end{enumerate}

\section{Evidence and source trail}

The detailed theorem-by-theorem ledger, including the named no-retry
constraints and the input required to revive each route, is
\href{https://github.com/gakonst/nanocodex-erdos700/blob/main/docs/part-iii-exploration-map.md}
{\texttt{docs/part-iii-exploration-map.md}}.  The compact specialist handoff is
\href{https://github.com/gakonst/nanocodex-erdos700/blob/main/docs/erdos700-iii-bottleneck-brief.md}
{\texttt{docs/erdos700-iii-bottleneck-brief.md}}.  Persisted-run and session
coverage audits are linked from the repository research map.

The primary historical source is P.~Erd\H{o}s and G.~Szekeres,
``Some number theoretic problems on binomial coefficients,''
\emph{Australian Mathematical Society Gazette} 5 (1978), 97--99,
\href{https://www.renyi.hu/~p_erdos/1978-46.pdf}{archival PDF}.

\end{document}
